\documentclass[../../main2.tex]{subfiles}

\begin{document}

\chapter{通道四：法律（博弈结果的定责编码）}
\label{ch:law}

\begin{motivation}
	上一章遗留的问题：政治裁决不可标价的冲突，可每次重吵一遍成本太高。社会需要把「吵过的结果」\textbf{写死}。\\
	\textbf{核心动机}：看法律作为第四个齿轮怎么咬合前三者，并给 Part IV 收尾——四条通道不是四个领域，是同一动力学按约束嵌套的四张脸。\\
	\textbf{角度提醒}：Part III 里法律是\textbf{方法的证明}（法庭逼出了分账手艺）；本章里法律是\textbf{系统的一个齿轮}（它在回路里怎么转）。两次出现，两个镜头，不是重复。
\end{motivation}

\begin{logicmap}
\begin{tikzpicture}[node distance=0.85cm, every node/.style={draw, rectangle, rounded corners, align=center, font=\small}]
\node (q) {裁决怎么写死？};
\node (s) [below=of q] {规则 = 定责的条件概率编码};
\node (sd) [below=of s] {与三齿轮的咬合};
\node (h) [below=of sd] {法律也会反噬：条文的意外后果};
\node (c) [below=of h] {收束：四通道的嵌套};
\draw[-{Stealth}] (q) -- (s);
\draw[-{Stealth}] (s) -- (sd);
\draw[-{Stealth}] (sd) -- (h);
\draw[-{Stealth}] (h) -- (c);
\end{tikzpicture}
\end{logicmap}

\section{规则：把定责写死}

\begin{readingnote}
	你有没有想过，交通法为什么精确到「红灯停」，而不是「红灯时视情况通过」？\\
	模糊的裁决要吵，\textbf{写死的规则}不吵——社会在买「不吵」。
\end{readingnote}

第 8 章看过法庭怎么\textbf{算}责任。这一节看社会怎么\textbf{存}责任：把反复吵出的结果编码成条文。用本书的语言，规则是\textbf{定责的条件概率编码}——「若行为 A，则损害责任概率几成」被提前写死，省得每次开庭。

\begin{tcolorbox}[colback=green!5, title=\textit{生活类比}]
	法律像游戏的规则手册。手册不告诉你这局牌怎么打（那是预测），只告诉你「出了这张牌算几分、违规怎么罚」（那是定责）。打牌的人按手册预期彼此的行为——预期一稳，交易成本就降。
\end{tcolorbox}

三个层次的编码，各配一例：

\begin{itemize}
	\item \textbf{行为定责}：闯红灯全责（推定的 $C$ 接近 1）——写死，不查动机。
	\item \textbf{比例定责}：交通事故的主次责七三开——第 9 章的分账，条文化。
	\item \textbf{程序定责}：证据不足时「疑罪从无」——承认 $C$ 估不出来时，宁可漏判不可错杀。这是\textbf{对不确定性的编码}，法学版的诚实标注。
\end{itemize}

\paragraph*{逻辑衔接}
	规则是定责编码。那它跟另外三个齿轮怎么咬？下一节看回路。

\section{与三齿轮的咬合}

法律不是终点，是回路里的齿轮。三根轴，逐个看：

\begin{itemize}
	\item \textbf{政治 $\rightarrow$ 法律}：立法是政治裁决的固化。税法、劳动法，都是吵赢了的一方把结果写死。\textbf{法律 $\rightarrow$ 政治}：程序给下一轮博弈定规则——选举法本身就是「怎么争」的规则。
	\item \textbf{经济 $\rightarrow$ 法律}：产权法、合同法是交易的基础设施。商路走到哪，商法跟到哪（第 12 章丝绸之路的罗马商法）。\textbf{法律 $\rightarrow$ 经济}：一纸产权界定，能改变全部激励——圈地、土改、专利法，都是法律直接改写经济参数。
	\item \textbf{模因 $\rightarrow$ 法律}：判例是「被赋予强制力的模因」——一个判决能训诫一代人。「昆山反杀案」改写正当防卫的社会认知，比一百篇普法文章管用。\textbf{法律 $\rightarrow$ 模因}：禁令也会制造模因——禁书榜永远是畅销榜。
\end{itemize}

\begin{questionbox}
	\textit{现在你可能想反驳：「法律条文写得那么死，哪是什么『概率编码』？它就是命令。」}\\
	\textbf{命令是壳，预期是瓤。}条文真正的作用，是让所有人\textbf{提前知道}「做了 A 会有什么下场」——这就是一个「下场的条件概率」被公告天下。死刑条文的威慑力，从来不在执行那天，在每个人动手前的那一下心跳。
\end{questionbox}

\paragraph*{逻辑衔接}
	咬合看完了。但齿轮也会打滑——写死的规则，会被下一轮的动机钻空子。看两个反噬，再收束 Part IV。

\section{条文的意外后果}

\begin{readingnote}
	你有没有想过，为什么每条新规出来，第二天就有人研究「怎么绕」？\\
	规则改变动机，新动机改造规则——这是第 16 章「循环」的提前放映。
\end{readingnote}

三个反噬现场：

\begin{itemize}
	\item \textbf{古}：王莽的「王田制」——土地国有、禁止买卖，条文完美，执行全面变形，豪强换个名目继续兼并。写死的规则敌不过活着的利益。
	\item \textbf{古}：明代「海禁」——禁出海，于是走私与海盗合流（倭寇里大半是沿海商人）。规则制造了自己的反面。
	\item \textbf{今}：学区房——「就近入学」的规则本意公平，配上稀缺的名校，催生了房价的军备竞赛。\textbf{规则 + 稀缺 = 套利空间}。
\end{itemize}

\lowfreq{（规律雏形：规则塑造的预期越稳定，钻空子的动机越强——因为套利收益变可靠了。这是猜想级的机制，没有定量验证。）}

\paragraph*{逻辑衔接}
	四个齿轮全部拆完。现在到了 Part IV 的收尾时刻——把四条通道的关系钉成一句可检验的话。这不是修辞总结，是推导结果。

\section{收束：四通道是一个嵌套}

从第 10 章的生成树往回读，嵌套关系是被\textbf{推出来}的，不是被声明的：

\begin{quote}
	\textbf{观念}（模因）$\supset$ \textbf{经济}（观念 + 稀缺）$\supset$ \textbf{政治}（经济 + 不可标价偏好）\\
	\textbf{法律}：为三者的博弈结果提供定责的条件概率编码。
\end{quote}

含义逐条落地：

\begin{itemize}
	\item 一切经济行为都是观念传播（买什么，先信什么）——但观念传播不都是经济（谣言不标价）。\textbf{经济是观念空间里被稀缺约束切出的子集}。
	\item 一切政治冲突都在经济背景里烧（名额、预算、土地）——但经济交换不都是政治（买白菜没人动员）。\textbf{政治是交换空间里被不可标价偏好切出的子集}。
	\item 法律不在三层\textbf{之内}，它是横跨三层的\textbf{编码器}：给观念战定诽谤边界、给交换定产权、给冲突定程序。
\end{itemize}

为什么说「同一张脸」？因为四条通道跑的是同一套动力学：\textbf{认同 + 资源 + 传播}。模因只跑传播；经济加上资源的守恒；政治加上认同的裁决；法律给三者的平衡态拍快照、写存档。四个领域教科书各写一本，动力学只有一本。

\begin{warningbox}
	以上是\textbf{统一视角的直觉呈现}。三个命题（经济是稀缺化的模因、政治是人情化的经济、法律是博弈结果的条件概率编码）的严格推导，见数学附录与待发表的《介观统一论》。\lowfreq{本书只描述，不假装在证明。}
\end{warningbox}

\paragraph*{逻辑衔接}
	Part IV 闭环：微观动机沿通道长成宏观结构，四通道嵌套成一张脸。手里现在有完整的「因素清单」了。最后一问：拿着这张贡献表，怎么说出\textbf{还没发生}的事？第 15 章：从贡献表到未来走向。

\begin{thinkbox}
\begin{enumerate}
	\item 挑一条你天天打交道的规则（交通、校规、公司报销、平台条款）：它在编码「哪一类责任」？它咬合的是哪个齿轮——观念、资源，还是认同？谁最有动机钻它的空子？
	\item \textbf{反事实}：如果「疑罪从无」从一开始就是人类默认——没有过刑讯、没有过有罪推定——司法史会少哪些大案？社会会更安全还是更不安全？用「对不确定性的编码」推一遍。
	\item \textbf{反事实}：如果学区房规则改成「全市抽签入学」——同一所名校、同一群家长——房价的军备竞赛会熄火吗？还是换个赛道（补习、迁户口）？「规则 + 稀缺 = 套利空间」这条雏形规律，帮你预测一下。
\end{enumerate}
\end{thinkbox}

\end{document}
